
\section{Ablation Study on Sampling Strategy}
\label{sec:sampling_ablation}

\begin{table*}[h]
    \scriptsize
    \centering
    \captionsetup{font=scriptsize}
    \caption{Ablation Study on Calibration Sampling Strategies. We analyze the impact of different frame selection strategies for constructing the synthetic calibration set. The results demonstrate that our \textbf{Uniform} sampling effectively captures global statistics, outperforming temporally biased methods (Consecutive) and achieving comparable robustness to computation-heavy clustering methods (K-Means) across real-world benchmarks.}
    % \vspace{-.8em}
    \label{tab:quantmacvo:ablation_sample}
    \resizebox{0.8\linewidth}{!}{
    \begin{tabular}{l|c|c|cc|cc|cc|ccc}
        \toprule
        \multirow{2}{*}{\textbf{Sampling Strategy}} & \multirow{2}{*}{\textbf{Description}} &\multirow{2}{*}{\textbf{W/A Bits}} & \multicolumn{2}{c|}{\textbf{TartanAir (Sim)}} & \multicolumn{2}{c|}{\textbf{EuRoC (Real)}} & \multicolumn{2}{c|}{\textbf{KITTI (Real)}} & \multicolumn{2}{c}{\textbf{Avg.}}\\
        &  & & $r_{\mathrm{rel}}$ & $t_{\mathrm{rel}}$ & $r_{\mathrm{rel}}$ & $t_{\mathrm{rel}}$ & $r_{\mathrm{rel}}$ & $t_{\mathrm{rel}}$ & $r_{\mathrm{rel}}$ & $t_{\mathrm{rel}}$ \\
        \midrule
        \rowcolor{gray!5} \multicolumn{2}{l|}{Baseline}&FP32 / FP32 & .0790 & .2670 & .0028 & .0468 & .0164 & .0360 & .0351 & .1246 \\
        \midrule
        \multicolumn{10}{l}{\textit{Quantized (W8A8) with different calibration sampling:}} \\
        Consecutive & \footnotesize{Sequential frames (Temporal Bias)} &INT8 / INT8 & .1250 & .4520 & .0045 & .0850 & .0210 & .0520 & .0541 & .2098 \\
        Random & \footnotesize{Random selection} &INT8 / INT8 & .0811 & .2938 & .0029 & .0487 & .0162 &.0368  & .0359 & .1352 \\
        K-Means & \footnotesize{Feature clustering (Diversity)} &INT8 / INT8 & .0730 & .2676 & .0028 & .0475 & .0161 & .0363 & .0328 & .1250 \\
        \midrule
        \rowcolor{gray!20} Uniform & \footnotesize{Uniform sampling (same interval)} &INT8 / INT8 & .0720 & .2630 & .0029 & .0490 & .0161 & .0369 & .0325 & .1240\\
        \bottomrule
    \end{tabular}
    }
    \label{tab:quantmacvo:calib_sample}
\end{table*}

To determine the suitable strategy for constructing the global calibration set $\mathcal{D}_{calib}$, we evaluate four sampling protocols: \textit{Consecutive} (sequential frames), \textit{Random}, \textit{K-Means} (feature clustering for diversity), and Uniform sampling. As presented in \tref{tab:quantmacvo:ablation_sample}, the sampling strategy fundamentally determines the quantization robustness. The \textit{Consecutive} strategy suffers from severe degradation, yielding the highest average relative translation error. This confirms that calibration sets dominated by temporally correlated frames fail to capture the long-tail distribution of the domain, leading to overfitting on local statistics. While \textit{K-Means} clustering effectively maximizes data diversity and recovers performance, it incurs significant computational overhead for feature extraction and clustering. In contrast, Uniform strategy achieves comparable overall performance, slightly outperforming K-Means.
    
This result empirically validates that simple uniform sampling over diverse synthetic sequences is sufficient to approximate the global ``quantization envelope.'' It breaks temporal correlations and covers extreme motion dynamics without the computational cost of complex clustering, making it the most practical choice for efficient Sim-to-Real calibration.


\section{Selected Layer List}
\label{sec:selected_layers}

% Selected quantized operators by architectural component. Preamble: \usepackage{booktabs}
\begin{table}[h]
  \centering
  \caption{Selected layers for Hadamard rotation in the QuantMAC-VO front-end (depth 12; 28 layers are selected). All remaining layers are quantized with SmoothQuant.}

  \label{tab:quantmacvo:selected_layers}
  \resizebox{0.5\linewidth}{!}{
  \begin{tabular}{@{}r p{0.92\columnwidth}@{}}
  \toprule
  \multicolumn{2}{@{}l@{}}{\textbf{Context Encoder -- MLP}} \\
  \midrule
  19 & \ttfamily svt.blocks.1.1.mlp.fc1 \\
  \midrule
  \multicolumn{2}{@{}l@{}}{\textbf{Memory Encoder -- Feature encoder MLP}} \\
  \midrule
  20 & \ttfamily feat\_encoder.svt.blocks.0.0.mlp.fc1 \\
  21 & \ttfamily feat\_encoder.svt.blocks.0.0.mlp.fc2 \\
  \midrule
  \multicolumn{2}{@{}l@{}}{\textbf{Memory Encoder -- Patch embedding / attention context projection}} \\
  \midrule
  22 & \ttfamily patch\_embed.proj.4 \\
  23 & \ttfamily vertical\_encoder\_layers.0.local\_block.attn.context\_proj \\
  24 & \ttfamily vertical\_encoder\_layers.0.global\_block.attn.context\_proj \\
  25 & \ttfamily vertical\_encoder\_layers.1.local\_block.attn.context\_proj \\
  26 & \ttfamily vertical\_encoder\_layers.1.global\_block.attn.context\_proj \\
  27 & \ttfamily vertical\_encoder\_layers.2.local\_block.attn.context\_proj \\
  28 & \ttfamily vertical\_encoder\_layers.2.global\_block.attn.context\_proj \\
  \midrule

  \multicolumn{2}{@{}l@{}}{\textbf{Memory Decoder -- Covariance update \& GRU gating}} \\
  \midrule
  1  & \ttfamily cov\_update.cov\_head.conv2 \\
  5  & \ttfamily cov\_update.gru.convq2 \\
  7  & \ttfamily cov\_update.gru.convq1 \\
  9  & \ttfamily cov\_update.gru.convz1 \\
  10 & \ttfamily cov\_update.gru.convr1 \\
  11 & \ttfamily cov\_update.gru.convz2 \\
  12 & \ttfamily cov\_update.gru.convr2 \\
  \midrule
  \multicolumn{2}{@{}l@{}}{\textbf{Memory Decoder -- Encoder / Flow head / Aggregation}} \\
  \midrule
  2  & \ttfamily update\_block.flow\_head.conv2 \\
  3  & \ttfamily update\_block.encoder.conv \\
  13 & \ttfamily update\_block.encoder.convf2 \\
  8  & \ttfamily update\_block.aggregator.to\_v \\
  \midrule
  \multicolumn{2}{@{}l@{}}{\textbf{Memory Decoder -- GRU \& projection}} \\
  \midrule
  4  & \ttfamily update\_block.gru.convq2 \\
  6  & \ttfamily update\_block.gru.convq1 \\
  14 & \ttfamily update\_block.gru.convz1 \\
  15 & \ttfamily update\_block.gru.convr1 \\
  16 & \ttfamily update\_block.gru.convz2 \\
  17 & \ttfamily update\_block.gru.convr2 \\
  18 & \ttfamily proj \\
  \bottomrule
  \end{tabular}}
\end{table}


\section{PTQ Profiling Time}
\label{sec:ptq_profiling}

In the Methodology section, we highlight a key limitation of layer-wise sensitivity profiling: it substantially increases PTQ profiling time. Because layer-wise profiling quantizes and evaluates one layer at a time, it requires repeated forward passes that scale linearly with the number of layers. In contrast, our method estimates layer importance using only full-precision forward passes over the analysis set, resulting in much lower profiling cost. Our method completes profiling in under one minute and is 460$\times$ faster than the layer-wise baseline on an RTX 3090 Ti using the same analysis set (32 samples).

\begin{table}[htbp]
\centering
\captionsetup{font=small}
\caption{PTQ profiling time (min) comparison between layer-wise sensitivity method and ours.}
\begin{tabular}{lcc}
\toprule
 & Layer-wise & Ours \\
\midrule
Time (min) & 227.12 & 0.48 (460$\times$) \\
\bottomrule
\end{tabular}
\label{tab:quantmacvo:profiling_time}
\end{table}

